\documentclass[11pt,a4paper]{article}
\usepackage[T1]{fontenc}
\usepackage[utf8]{inputenc}
\usepackage[margin=2.2cm]{geometry}
\usepackage{amsmath,amssymb,booktabs,longtable,tabularx,array}
\usepackage[hidelinks]{hyperref}
\newcommand{\code}[1]{\texttt{\detokenize{#1}}}
\newcommand{\loc}[2]{\par\textbf{Location:} \code{#1}, lines #2.\par}
\newcommand{\finding}[4]{\subsection{#1: #4}\label{issue:#1}\textbf{Severity: #2. Evidence: #3.}\par}
\newcommand{\repair}{\par\textbf{Targeted repair:} }
\newcommand{\checklater}{\par\textbf{Acceptance check (not run):} }
\setlength{\parindent}{0pt}
\setlength{\parskip}{0.5em}
\setlength{\emergencystretch}{3em}
\title{CPD Implementation Audit: Ranked Issues and Repair Boundaries}
\author{}
\date{7 November 2026}
\begin{document}
\maketitle

\section{Scope, evidence and interpretation}

This is a static audit of revision
\code{c9d46ecbaea3500ee5cf1d3b8b278a9d780dd0d1}.
Line numbers refer to that checkout. No project code, tests, simulations, data
processing was executed. The document was subsequently compiled at the user's
request; compilation is not validation of CPD behavior. No pipeline code was modified.

The initial audit used the pipeline-level explanation in \code{docs/pipeline.tex}: patient-specific
unsupervised adaptation may use a held-out patient's available past; clinical
classifiers train on training patients; onset features must not use future readings;
gaps are defined on actual observations; comparisons must have explicit feature
bundles. Dedicated historical CPD notes are explanatory context, not a requirement
to restore an obsolete algorithm. Commit \code{8095fe3} deliberately replaced
whole-sequence HMM EM/Viterbi with fixed-parameter forward inference and introduced
the present whitening/calibration design.

\textbf{PDF verification, 7 November 2026.} The actual
\code{docs/pipeline.pdf} was subsequently extracted and checked directly.
It is the exported \emph{Asthma Prediction Protocol / Interactive Asthma Prediction
Report}, not a compiled copy of \code{pipeline.tex}. Its PDF metadata identifies
Chrome/Skia as the producer and a creation date of 4 November 2026. It contains two
very tall PDF pages; the section numbers below are more useful than page numbers.
The PDF is the system-design reference for the alignment table that follows.

\section{Explicit alignment with the system pipeline PDF}
\begin{longtable}{@{}p{0.19\textwidth}p{0.40\textwidth}p{0.34\textwidth}@{}}
\toprule PDF section & Required system behavior & Relevant audit findings \\
\midrule\endhead
1: objectives & Detection asks about day D; onset asks about the following
1--3 days. Features include day D, not subsequent observations. & C01, C05, C06,
C11, C19 address information availability. Full-day features themselves are allowed. \\
2: available measurements; 4: invariant variables & Every path gives the classifier
all specified measurements. The processing path changes HR/HRV features; temp,
steps and sleep remain available raw. & C03 is a direct implementation-contract
failure. C18 identifies branch-dependent support of nominally raw drift inputs. \\
3: signal processing; 3.7: phases & Compute each channel's derived signals once,
store them, then assemble paths without repeating work per classifier/fold. & C04
requires correct content dependencies, not moving all CPD fitting into every fold.
Patient-local past-only adaptation is compatible with shared cached artifacts. \\
3.1; 4: path matrix and path 5 flow & Smoothing supplies raw, baseline and residual;
CPD reads the path's designated signal and emits features. & C03, C18 and C24
check preservation, support and aggregation. Nightly raw-level drift is an added
implementation branch whose role should be explicit in each comparison. \\
3.2/3.5; A.2: lookahead & Whole-record smoother calibration is recognized as
retrospective; real-time settings must use past data. & C06 is a conditional
real-time limitation already acknowledged by the PDF, not a demand to discard
retrospective smoother experiments. C01 and C05 are additional CPD-local leaks. \\
4: invariant patients and split & Keep all 191 patients under the stated default
eligibility; unavailable method outputs are blank, not removed days. Use the same
patient split across paths. & C28 repairs must preserve blank-feature rows and
eligibility. Successful-only cohort restriction is not an acceptable repair. \\
4: tasks, windows and horizons & Keep tasks separate; window changes add history;
horizon changes narrow labels without moving the denominator. & This audit proposes
no target, window or eligibility change. CPD validity must not silently select
different scoreable days. \\
4/5: scoring & Use consistent day, episode and probability-band evaluation. & C09,
C14--C16, C20, C24 and C29 concern CPD event/exposure definitions and their
features/reports. CPD alarms are not clinical episode labels or probabilities. \\
A.3: gap rule & A break of at least 180 minutes starts a new segment. Channels
retain their own time grids. & C02 and C08 directly concern segmentation;
C07, C13, C19, C25 and C32 concern how detector memory/output respects those
boundaries. A common grid across channels is not proposed. \\
A.4: architectural rules & Reuse keyed computations and preserve manifests,
modularity and controlled comparisons. & C04 and C31 strengthen provenance;
targeted repairs should preserve stage separation and method discovery. \\
\bottomrule
\end{longtable}

\textbf{What the PDF does not specify.} It does not provide a complete mathematical
contract for whitening, null generation, HMM fitting, CUSUM recovery, analytic drift
calibration or every alarm lifecycle. Its references to four detectors also lag the
five-method implementation. Therefore R-labelled findings below are implementation
assumptions requiring validation or a design decision, not invented PDF requirements.
The five-method code is audited on its own stated behavior, without requiring a
return to an older four-detector design. Proposed streaming interfaces are optional
repair architecture, not a condition imposed by the PDF.

\textbf{Alignment is not certification.} This mapping establishes that the audit
targets the intended system. It does not establish correct detector mathematics,
useful clinical signal, or an acceptable false-alarm/delay trade-off. Those require
the regression and independent null/injection checks below, followed by the PDF's
held-out predictive evaluation on unchanged patients, days and feature bundles.

\section{Evidence and severity conventions}

\textbf{Evidence labels.} D = a control-flow, data-flow or algebraic defect visible
in source, with an explicit triggering condition; no claim about its observed
frequency. R = a statistical/model-design risk requiring validation or a policy
decision, rather than a proven coding error. V = a verification or observability
gap. Some findings carry more than one label. A proposed fixture is a specification
for later verification, not a test result. No finite static review guarantees that
all bugs have been found, and this report does not estimate clinical harm or metric
inflation from unavailable data.

\textbf{Severity ranks by potential damage to conclusions and processing.}
S1 (critical) can invalidate causal claims or silently change/reuse an experimental
input. S2 (major) can systematically corrupt a detector, null, alarm timing or
calibration. S3 (moderate) distorts interpretation or affects particular edge cases.
S4 (minor) affects robustness, reporting or maintenance. These are engineering
priorities, not measured clinical severity. Within a tier, fix order depends on the
branches actually used.

\section{Ranked register}
\small
\begin{longtable}{@{}p{0.08\textwidth}p{0.09\textwidth}p{0.08\textwidth}p{0.68\textwidth}@{}}
\toprule ID & Rank & Evidence & Issue \\
\midrule\endhead
C01 & S1 & D & Shuffle calibration reads beyond its activation timestamp \\
C02 & S1 & D & Raw filler rows hide observation gaps \\
C03 & S1 & D & Channel metadata and feature-coverage bookkeeping are wrong \\
C04 & S1 & D & CPD cache omits shared code and side-input contents \\
C05 & S1 & D & Future sleep availability changes earlier drift windows \\
C06 & S1 & D/R & Upstream lookahead remains in CPD inputs \\
C07 & S2 & D & Held HR output crosses resets and settling boundaries \\
C08 & S2 & D & Shuffle null replaces original chunks with calendar days \\
C09 & S2 & D/R & Calibration measures events but permits nearly continuous alarms \\
C10 & S2 & D & Missing scores can become alarms at nonpositive thresholds \\
C11 & S2 & D & KCPD backdates scheduled results and revises earlier outputs \\
C12 & S2 & D/R & Whitening carries latent states across changed coordinates \\
C13 & S2 & D/R & Kalman CUSUM carries evidence indefinitely across gaps \\
C14 & S2 & D & Consecutive nightly drift alarms are suppressed \\
C15 & S2 & D & Drift loses/delays nights and collapses simultaneous verdicts \\
C16 & S2 & R & Drift's analytic budget does not calibrate its actual process \\
C17 & S2 & D/V & STS accepts unreliable fits and has incomplete fallback handling \\
C18 & S2 & D/R & STS null support is conditioned on upstream output availability \\
C30 & S2 & V & Existing checks do not establish the advertised guarantees \\
C32 & S2 & D & Long nonwear resumes with an empty baseline but remains armed \\
C19 & S3 & D & Partial-bin closure depends on later input or end-of-file \\
C20 & S3 & D/R & Worn-time and sleep-coverage accounting use nominal counts \\
C21 & S3 & D/R & Tail fitting discards scale; HMM uses a fixed scale thereafter \\
C22 & S3 & D/R & Calibration exclusions/fallbacks can include abnormal or old history \\
C23 & S3 & D/R & Severity values are not comparable across detectors \\
C24 & S3 & D & Daily event counts and time labels disagree with CPD accounting \\
C25 & S3 & D/R & BOCPD gap semantics, phenotype and pruning need explicit bounds \\
C26 & S3 & R & STS synthetic fidelity and innovation-tail assignment are unverified \\
C27 & S3 & D/V & Invalid settings and partially malformed inputs are not rejected \\
C28 & S3 & D/V & A patient error can abort a cohort and leave a running artifact \\
C29 & S3 & D/V & Calibration reports can compare different inputs and exposure sets \\
C33 & S3 & D/V & Missing one state source silently discards the other \\
C31 & S4 & D/V & Provenance and output omit important diagnostic state \\
\bottomrule
\end{longtable}
\normalsize

\section{What is trained, fitted, updated or fixed}
\begin{tabularx}{\textwidth}{@{}p{0.19\textwidth}XX@{}}
\toprule Component & Current adaptation & Audit interpretation \\
\midrule
BOCPD & Normal--Gamma posterior per run hypothesis; fixed prior and hazard & Online Bayesian updates are learning; no batch fit is missing. \\
KCPD & Window-specific bandwidth and permutation mean/SD & Local statistical estimation, not supervised training. \\
Kalman monitoring & Spike, cumulative sums, trailing pair statistic & Its fitted state-space noise model is upstream. \\
HMM & Forward probabilities; fixed dwell times and variance ratio & No EM by design; tail degree of freedom comes from patient burn-in. \\
Nightly drift & Rolling median/MAD and CUSUM & A fitted personal baseline plus an analytic threshold. \\
Whitening & Hour/state profiles, OU parameters, initial tail fit & Fitted on patient history at arming and later refits. \\
Shuffle & Thresholds fitted on resampled patient history & Shared unsupervised calibration affects all simulated streams. \\
STS & Gaussian state-space likelihood fit, separate tail fit, simulated thresholds & Genuine batch fitting at each calibration; must use past data only. \\
\bottomrule
\end{tabularx}

None of these CPD components trains on exacerbation labels. This is consistent with
the pipeline. Whether HMM dwell times/emission parameters should instead be learned
is a model-selection question. Reintroducing full-record EM or Viterbi would restore
lookahead rather than repair it. Any supervised choice of detector settings must
respect training/validation patient boundaries.

\section{Critical issues}

\finding{C01}{S1}{D}{Shuffle calibration exceeds its historical cutoff}
\loc{wearable/steps/detect.py}{222--237}
\loc{wearable/cpd/_whiten.py}{316--345}
\code{_thresholds} selects dates from indices before \code{p0}, but passes full-length
arrays to \code{twin}. Its pools select \emph{all} rows on each selected date.
Arming midway through a day therefore includes later readings on that date.
The ordinary later recalibration points are first-of-day indices, so their current
date is not normally present in the prefix: the direct future-in-day counterexample
is initial mid-day arming, not every weekly recalibration. Even without future
rows, selecting a boundary date can include observations before the trailing-window
lower bound. STS's direct slice is timestamp-bounded; its shuffle fallback inherits
the shuffle defect. This can affect thresholds applied before their source data existed.
\repair Pass an explicitly bounded history object, with \code{stamp < t_cal} and
the intended lower bound, into every null generator. Dates must never be a substitute
for timestamp bounds.
\checklater Arm at 12:00, change only observations after 12:00, and require identical
thresholds and earlier outputs. Repeat with STS failure forcing shuffle fallback.

\finding{C02}{S1}{D}{Raw-path segmentation is computed before removing fillers}
\loc{wearable/steps/detect.py}{86--90 and 286--297}
\loc{wearable/steps/ingest.py}{60--76}
Ingest inserts NaN filler rows every ten minutes across long gaps. Detection assigns
chunks on adjacent \emph{rows}, then removes nonfinite signal rows. Thus a four-hour
observation gap becomes a series of ten-minute row gaps and no H8 boundary.
Whitening misses its hard reset; KCPD's no-cross-boundary rule has no correct IDs to
use. Exact 180-minute gaps can split while longer, filled gaps do not. Existing
partially populated chunk columns are also accepted without validating every observed
row. The separate chunk implementation defeats the shared segmentation contract.
\repair Construct/validate chunk IDs on the relevant actual observations with
\code{wearable.chunking.assign_chunks}; map them to presentation rows afterwards.
Specify whether missing transformed outputs retain raw-signal chunks.
\checklater Identical observed series with and without NaN fillers must produce
identical observed-row chunks and detections, including 179/180/181-minute gaps.

\finding{C03}{S1}{D}{Artifact coverage changes the available model variables}
\loc{wearable/steps/detect.py}{335--344 and 513--517}
\loc{wearable/steps/join.py}{195--230}
Detection does not record \code{stats.signal}; join defaults to HRV when deciding
which channel's raw block is covered. An HR-only detector can suppress raw HRV and
duplicate raw HR. With the default two raw+CPD channels, HR is duplicated even when
HRV happens to be present in its own detector table. Separately, the detected table
retains only the chosen input column, not necessarily raw measurements. Merely adding
\code{stats.signal} will therefore not repair all cases: a standalone residual,
smoothed, surprise or tail detection artifact can suppress a raw block it does not
contain. Standard experiment routes also pass smoother/forecast artifacts, which
mitigates this second case but not the incorrect detector metadata.
\repair Declare both signal identity and actual raw-feature coverage. Prefer joining
a canonical raw block independently, with transformed features explicitly additive.
Do not infer contents from channel identity alone.
\checklater Enumerate actual feature columns for HR-only, HRV-only and two-channel
raw+CPD, residual+CPD and surprise+CPD, both standalone and through experiment recipes.
All intended raw variables must appear once, on the same day support.

\finding{C04}{S1}{D}{CPD computation and dependencies are under-keyed}
\loc{wearable/steps/detect.py}{259--267 and 455--470}
\loc{wearable/store.py}{132--140}
Only detector-module hashes enter the CPD code hash. Changes to
\code{steps/detect.py}, \code{_whiten.py}, \code{_sts.py}, or shared segmentation can
reuse old detections. Raw input is keyed by the directory string rather than its
contents; sleep/activity side files are read outside the declared upstream artifact.
Their addition, correction or replacement changes results without changing the key.
Recorded Git revision is provenance, not invalidation. Existing downstream artifacts
can consequently survive a scientifically material CPD correction.
\repair Version/hash the complete CPD implementation and content-address the actual
channel and state inputs. Introduce new logic keys when fixing these issues; rebuild
dependent joins/models/evaluations with a documented invalidation boundary.
\checklater Changing each shared helper or one state-file value must change the
appropriate CPD key; unchanged content must reuse it. No raw data need be committed.

\finding{C05}{S1}{D}{Future sleep flags select the definition of past nights}
\loc{wearable/cpd/drift.py}{71--83}
\code{has_sleep} is computed across the entire supplied state array. A prefix with
no sleep uses 00:00--06:00 raw readings. Appending the first sleep flag switches all
earlier nights to sleep-only readings in the 18:00--12:00 window, potentially deleting
those nights and changing baseline, alarms and magnitude. This is a CPD-internal
causality failure, independent of smoother lookahead.
\repair Make availability/fallback decisions at each night's decision time using
only state observations available then, with a declared minimum coverage policy.
\checklater Append the first sleep observation several days after arming; all
previous verdicts and their baselines must remain unchanged.

\finding{C06}{S1, conditional on input}{D/R}{Causal detection cannot repair leaked inputs}
\loc{wearable/experiment.py}{270--287}
\loc{wearable/steps/detect.py}{374--394}
The default smoothing recipe is \code{gammadglm}; the existing causality audit
documents whole-frame parameter fitting in nine smoothers. Detection accepts those
artifacts without a causal-status check. Their early residuals, and histories used
to fit whitening/thresholds, can already depend on later observations. This is an
upstream issue at the CPD boundary, not evidence that every detector recurrence is
noncausal. It particularly affects onset interpretation.
\repair Record input causal status and calibration interval in manifests; distinguish
retrospective comparisons from prospective-valid branches. Use causal initialization
or a fixed past calibration for the smoother before certifying the chain.
\checklater Full-chain prefix consistency must include recomputing the upstream
transform, rather than truncating a residual computed on a full recording.

\section{Major algorithm and calibration issues}

\finding{C07}{S2}{D}{HR output holding leaks old flags across a reset}
\loc{wearable/cpd/_whiten.py}{302--313}
\loc{wearable/steps/detect.py}{324--359}
The holder maps a reading to the last bin's final-row index, regardless of chunk or
elapsed gap. The first readings after a four-hour gap, before their new bin's last
row, therefore hold the \emph{previous chunk's} output, including its armed flag.
An old alarm can appear during the new settling interval. A shorter re-entry has the
same issue. At midnight, old held outputs can also be attributed to the new day.
This is not fixed by correctly segmenting the binned series.
\repair Apply current-reading chunk/re-entry/settling eligibility after mapping, and
allow holding only within an explicit valid-time interval. Never borrow an armed
state from another chunk. Keep verdict impulses distinct from held states.
\checklater End a chunk armed and flagged, then resume after a long gap with at least
two observations in the first bin; no resumed reading may inherit the old alarm.

\finding{C08}{S2}{D}{Shuffle null destroys internal resets and adds artificial ones}
\loc{wearable/cpd/_whiten.py}{328--345}
The supplied \code{chunk} array is unused: every sampled calendar day becomes one
chunk. Within-day H8 boundaries vanish; midnight becomes an H8 boundary even for
continuous wear. Noise filtering, settle exposure and KCPD validity thus differ
between null and real data. Repeating days also introduces joins and level transitions
that are not evidence of a physiological change in the original data.
\repair Preserve internal chunk boundaries and explicitly define/reset at synthetic
block joins. Calibrate with the same reset and availability rules as real inference;
document which original continuity is intentionally destroyed.
\checklater Map sampled days back to original rows and verify every internal reset,
all synthetic joins, settling masks and eligible KCPD windows.

\finding{C09}{S2}{D/R}{Event-rate calibration admits an always-on detector}
\loc{wearable/cpd/_whiten.py}{349--370}
\loc{wearable/steps/detect.py}{241--250}
Threshold choice constrains rising-edge counts only. A threshold producing one long
alarm throughout 60 worn days can meet a budget even though virtually all readings
are flagged. The ``at every higher threshold'' envelope helps reject low thresholds
when intermediate thresholds produce many events, but a flat/degenerate score has
no such protection. For a constant finite score, the minimum candidate can be
accepted whenever one event fits the budget. Also, each statistic is calibrated in
isolation; the final detector and ensemble are not replayed and checked on the null.
\repair Specify an event lifecycle, occupancy/duration constraint and admissible
threshold domain. Record both null event rate and occupancy, and replay final
detector/ensemble outputs rather than equating a low onset count with useful alarms.
\checklater Constant scores, sustained plateaus, isolated spikes and correlated
streams must not yield near-continuous accepted alarms without an explicit state
interpretation. Verify final OR behavior, not just per-stream thresholds.

\finding{C10}{S2}{D}{Unavailable statistics become alarms at zero/negative thresholds}
\loc{wearable/cpd/kcpd.py}{123--129}
\loc{wearable/cpd/kalman.py}{85--91}
KCPD and Kalman convert NaN scores to zero before thresholding. Calibration retains
missingness as non-hits, so inference has a different rule. If the learned threshold
is zero or negative, invalid KCPD windows can alarm. Constant zero innovations give
zero MMD scores; C09 can accept a zero threshold, so insufficient-coverage intervals
then become positive too. Dividing severity by zero also yields NaN/inf. Kalman has
the analogous problem for zero-valued streams. Outer arming does not establish that
each individual detector statistic is available.
\repair Require finite scores in every hit mask. Represent unavailable separately
from normal; make severity defined for all allowed thresholds.
\checklater A missing statistic can never alarm for any valid learned threshold;
exercise zero/negative thresholds and degenerate nulls explicitly.

\finding{C11}{S2}{D}{KCPD backdates evaluations and can overwrite earlier readings}
\loc{wearable/cpd/kcpd.py}{79--111}
For evaluation time $T$, the code writes to the last reading at or before $T$.
When $T$ lies between observations, that reading precedes availability of the
scheduled result. Several grid evaluations in a gap can map to the same index;
\code{last_val} is replaced and later written back onto it. A longer recording
extends the evaluation grid, potentially changing an earlier reading without a new
observation near it. The arithmetic uses data no later than $T$, but the output's
timestamp is still wrong.
\repair Keep evaluation timestamps as first-class output, or evaluate at actual
arrival times and publish only at/after availability. Do not map a later clock event
backwards. Reset/expire held statistics at declared gaps.
\checklater Use jittered arrivals, a long silent interval and many cuts inside that
interval; already-published scores must be immutable.

\finding{C12}{S2}{D/R}{Whitening refits leave the latent state in old coordinates}
\loc{wearable/cpd/_whiten.py}{203--228}
Hourly centers/scales and state corrections are recomputed at each refit, changing
the normalized observation coordinate. The OU filter carries its old normalized
\code{x} and \code{P} into that new coordinate unchanged. A center shift of one old
SD does not magically leave the old latent mean valid. New OU amplitude/timescale
also changes the interpretation of covariance. Midnight innovations can therefore
reflect a parameter-update artifact, not a new physiological change.
\repair Transform latent mean/covariance into the new normalization where possible,
or causally replay a bounded history under the new parameters before publishing
new innovations. A reset is another policy but must have explicit settling behavior.
\checklater Force a known normalization change with unchanged raw physiology and
compare against a correctly transformed/replayed filter; no artificial spike.

\finding{C13}{S2}{D/R}{Kalman drift evidence ignores elapsed time and resets}
\loc{wearable/cpd/kalman.py}{49--55}
The CUSUM simply adds each finite $z$; it neither uses elapsed time nor chunk IDs,
and does not restart after an alarm. A large accumulated value can persist across
hours/days of missingness, then combine with unrelated later evidence. There is no
clock-time decay or reset corresponding to the orchestration comment that detectors
handle elapsed time. The chosen step changes the evidence per hour. No-reset is
an intentional possible monitoring policy, but the gap/duration semantics here are
not equivalent to an H8 state flush or independent alarm episodes.
\repair Decide whether this is cumulative regime evidence or repeated change
detection; implement elapsed-time weighting/decay and explicit gap/rearm policy,
then recalibrate that same recurrence. Do not reset only at inference.
\checklater Separate two subthreshold excursions by a long gap; check whether their
combination is permitted by the chosen policy. Test recovery after a large alarm.

\finding{C14}{S2}{D}{Consecutive nightly alarms collapse into one}
\loc{wearable/cpd/drift.py}{126--137 and 148--154}
The CUSUM restarts after a verdict, but its above-threshold score remains held.
If the next night independently alarms, the held Boolean never falls; rising-edge
extraction deletes that second verdict. This contradicts once-per-verdict output.
\repair Return explicit verdict indices/IDs and one impulse for each alarmed night;
keep held magnitude separate from event identity.
\checklater Two and three consecutive alarmed nights must produce two and three
distinct verdict events, even with no below-threshold score between them.

\finding{C15}{S2}{D}{Night enumeration and delayed verdict mapping lose information}
\loc{wearable/cpd/drift.py}{76--91 and 112--137}
Only calendar dates containing readings are enumerated. An evening's recorded sleep
can yield a valid window ending on a date with no readings, but that night is never
examined. Other nights are delayed to the first armed reading, potentially days
later. Several nights can map to one index, where the held-score writer overwrites
earlier verdicts. Their CUSUM/baseline updates still happened internally, so emitted
output no longer describes the sequence processed. No maximum verdict delay exists.
\repair Generate candidate night-end dates from the window support; keep separate
night time, availability time and publication time. Expire stale verdicts or retain
them as explicit events rather than silently overwriting one row.
\checklater Recorded evening sleep followed by a no-reading next day; multiple
nights completing before the next armed reading; recovery after extended nonwear.

\finding{C16}{S2}{R}{Drift's analytic threshold does not calibrate its actual statistic}
\loc{wearable/cpd/drift.py}{50--59 and 114--132}
The analytic approximation assumes standardized nightly observations of a specified
no-change process. The actual $u_D$ uses a selected rolling median/MAD estimated
from as few as three nights, a one-device-unit floor, exclusion of alarmed nights
and missing-night skips. Its distribution need not be independent $N(0,1)$.
One eligible night is not one 24-hour worn day. Consequently assigning this
approximation one-fifth of the ensemble's worn-day budget does not verify its rate.
The analytic detector is exempt from empirical null-rate reporting.
\repair Calibrate the entire nightly algorithm on realistic past-only nulls, or
state a separate per-eligible-night sensitivity target and validate the approximation.
\checklater Vary baseline length, missing nights, sleep coverage, heavy tails and
state availability; measure false verdicts per eligible night and per worn day.

\finding{C17}{S2}{D/V}{STS fit health is unchecked and fallback is incomplete}
\loc{wearable/cpd/_sts.py}{126--139 and 178--191}
\loc{wearable/steps/detect.py}{155--187}
The likelihood fit suppresses warnings and accepts returned parameters without
checking optimizer convergence, finite values or useful variance estimates.
The exception guard covers the main fit but not the later Student-t fit or
simulation. Those failures escape instead of returning the documented shuffle
fallback. A finite but unusable fit can generate degenerate nulls and bad thresholds.
The shared initial tail fit in \code{_whiten.py:221--224} is likewise unguarded.
\repair Validate fits and simulations explicitly; guard the whole calibration
operation with narrowly classified errors; record fallback reason and fit diagnostics.
Do not silently convert a numerical failure into a normal detector result.
\checklater Nonconverged optimizer, constant series, near-zero variances, invalid
tail fits and simulation exceptions must produce declared fallback/failure status.

\finding{C18}{S2}{D/R}{STS and nightly drift are restricted by transformed availability}
\loc{wearable/steps/detect.py}{288--308 and 151--154}
All data are first filtered to finite values of the \emph{chosen detector column}.
Only then are raw measurements saved for STS and drift. Raw readings in short chunks
where a smoother emits NaN, or before a forecast becomes available, are absent.
Thus nightly drift, advertised as a raw-channel detector, changes its raw observation
support across branches. STS fits and simulates that restricted mask; raw observations
that would initialize a real smoother can be missing from its synthetic run.
This is more than choosing residual versus raw as a score input.
\repair Carry the full raw/state stream separately from transformed eligibility.
Define which detector needs which support; generate raw STS histories first, then
apply the actual transform and eligibility policy once.
\checklater Give two transforms different NaN masks over identical raw observations;
raw nightly baselines should agree unless the experiment explicitly varies support.

\finding{C32}{S2}{D}{A long absence leaves an armed detector with no fitted baseline}
\loc{wearable/cpd/_whiten.py}{194--217 and 235--236}
\loc{wearable/cpd/_whiten.py}{64--69 and 126--127}
Once arming has occurred it is never revoked. After an absence longer than
\code{baseline_days}, the historical slice at the next refit can be empty. The
fewer-than-ten fallback selects the same empty slice; it does not recover a valid
baseline. Robust fitting then returns center 0 and scale 1, with no OU correlation.
For a raw HR input near 80, this yields innovations near 80 rather than a meaningful
patient-standardized score. The patient becomes eligible again after only the short
settling interval, despite having no current fitted baseline. Earlier detector
memory remains in use. Threshold recalibration can resample arbitrarily old history
through C22. This differs from retaining state across a small gap.
\repair Require sufficient eligible fitting history after each material absence.
Freeze an explicitly valid historical fit with a staleness policy, or re-enter
calibration/unarmed status; never substitute generic center/scale while reporting
ordinary armed inference.
\checklater An armed raw-channel patient disappears for more than 14 days and
returns with unchanged physiology. Status and parameters must follow the declared
re-entry policy and must not produce default-zero-centered pseudo-anomalies.

\section{Moderate issues and model assumptions}

\finding{C19}{S3}{D}{A bin's closure time is inferred retrospectively}
\loc{wearable/cpd/_whiten.py}{302--313}
\loc{wearable/cpd/kcpd.py}{50--59}
Bins are timestamped at their final observed row, including the last incomplete
bin at end-of-file. That row is not known to be the last until the bin closes or a
later arrival establishes it. Truncating within a bin publishes a different partial
median than extending the input. Even a completed sparse bin's last arrival can
precede its known clock closure. This is an availability-time defect, not a centered
future median at every row. The selftest explicitly excludes the final bin near cuts.
\repair Emit only closed bins at a declared clock boundary, or label partial outputs
as provisional and exclude them from causal claims. Apply the same rule to KCPD's
optional second binning layer.
\checklater Prefix cuts at every offset in a five-minute bin, including sparse bins
and a closure with no new observation.

\finding{C20}{S3}{D/R}{Nominal sample counts are presented as actual time coverage}
\loc{wearable/cpd/_whiten.py}{185--197 and 331--346}
\loc{wearable/cpd/drift.py}{84--85}
\loc{scripts/cpd_calibration_compare.py}{55--64}
Arming uses $n\,\Delta t_{nominal}$, not measured worn duration. One HR observation in
a five-minute bin counts five worn minutes after binning; an irregular HRV sample
counts ten minutes regardless of actual support. Sleep coverage similarly multiplies
reading count by nominal step. Reported exposure on mapped original HR rows instead
uses one minute per row, so calibration and reported rates can have different
denominators under sparse wear. This proxy may be intentional but is not literal
24-hour wear, and duplicates can inflate it.
\repair Declare effective-sample exposure versus actual coverage; use one bounded
duration/coverage rule consistently for arming, null length, sleep eligibility and
reported rates. Never bridge arbitrary gaps by assigning their whole duration.
\checklater Equal sample counts with different occupancy, sparse HR bins and
duplicate timestamps; exposure definitions must be documented and consistent.

\finding{C21}{S3}{D/R}{Tail fitting and HMM emission scale are inconsistent}
\loc{wearable/cpd/_whiten.py}{218--224}
\loc{wearable/cpd/hmm.py}{38--47}
The Student-t fit estimates both degree of freedom and scale with location fixed,
but stores only degree of freedom. The HMM then assumes scale 1 in state 0 and
scale $\sqrt{v_1}$ in state 1. Unit variance and Student-t scale 1 are different:
$\operatorname{Var}(t_\nu)=\nu/(\nu-2)$. Robust normalization does not itself prove
either scale or variance is one. The initial degree of freedom is then frozen while
other noise parameters refit. Threshold calibration may compensate for some mismatch
but does not turn posterior state probability into a calibrated physiological risk.
\repair Choose and fit a consistent emission scale convention; retain fit scale and
state whether tails are frozen or refitted. Validate adaptation rather than adding
unrestricted full-record HMM training. Rename the stale ``Gaussian HMM'' label.
\checklater Known t-distributed innovations with multiple scales and tail weights;
verify likelihood ratios and calm-state behavior under the chosen convention.

\finding{C22}{S3}{D/R}{Baseline exclusion can silently abandon the healthy-history rule}
\loc{wearable/cpd/_whiten.py}{205--234}
\loc{wearable/steps/detect.py}{224--235}
If abnormal-day removal leaves fewer than ten rows, whitening reintroduces the
unfiltered window. If every candidate shuffle date is excluded, calibration uses
\emph{all} dates before \code{p0}, bypassing both exclusion and the nominal trailing
window. Slow deterioration can also remain below the abnormal-day criterion and be
absorbed by normal refitting. The final abnormal-date set is passed to calibration
without an as-of ledger. Normally past dates are already complete, so this alone is
not proof of broad future leakage; initial partial-day classification deserves an
explicit boundary check.
\repair Make fallback choices explicit and recorded; consider freezing the last
valid baseline rather than silently treating abnormal data as normal. Supply an
as-of exclusion ledger and bounded history to every recalibration.
\checklater Entire trailing window abnormal, fewer than ten healthy rows, persistent
subthreshold drift, and arming mid-calendar-day.

\finding{C23}{S3}{D/R}{Ensemble severity selects unlike quantities}
\loc{wearable/steps/detect.py}{94--117}
\loc{wearable/cpd/bocpd.py}{110--112}
\loc{wearable/cpd/hmm.py}{77--83}
\loc{wearable/cpd/kalman.py}{85--91}
\loc{wearable/cpd/kcpd.py}{123--129}
BOCPD/HMM use posterior probabilities; other methods use score divided by twice a
threshold. These have no common severity meaning. KCPD thresholds may be negative,
making that normalization misleading. Kalman's severity is the maximum over all
streams, while phenotype/magnitude use fixed spike-over-ACF-over-drift precedence;
its winning severity can describe another stream. NaN severity removes a firing
method's phenotype claim while its anomaly still contributes to the OR.
\repair Preserve all stream scores and select a representative using a declared
priority or comparable null-tail measure. Keep winning stream, phenotype and
magnitude coupled. Do not interpret primary phenotype as a diagnosed cause.
\checklater Simultaneous stream/detector hits, ties, negative/zero thresholds and
nonfinite severities must have deterministic, semantically consistent output.

\finding{C24}{S3}{D}{Event and duration definitions differ across pipeline stages}
\loc{wearable/steps/detect.py}{501--517}
\loc{wearable/steps/join.py}{166--172}
\loc{wearable/cpd/_whiten.py}{285--288}
Detection counts runs after deleting unarmed readings, intentionally merging across
settling windows; it also merges across long gaps if flags remain high. Join counts
on the full rows separately within each day: filler/unarmed zeros split events and
midnight starts another daily bout. Columns ending in \code{_minutes} sum flagged
rows rather than durations. HR bin holding expands a one-bin verdict onto several
rows, so drift is no longer a single-row impulse on the original grid. These are
different estimands unless explicitly named.
\repair Define event IDs and exposure/duration once, preserve them through mapping,
and choose whether a daily count means new event onsets or events active that day.
Rename sample counts or implement a documented bounded duration calculation.
\checklater One event across midnight, a gap and a settle interval; check global
counts, daily counts, impulses and durations against the declared definitions.

\finding{C25}{S3}{D/R}{BOCPD gap behavior and phenotype are weaker than their labels}
\loc{wearable/cpd/bocpd.py}{63--94 and 115--125}
The recurrence ignores chunk IDs and increases hazard with elapsed time across a
gap. A return reading can therefore favor a new segment even without observed
change; settling may suppress reporting but does not clear that evidence. This is
an intentional alternative to hard resets, so requires an explicit compatibility
decision with H8, not automatic replacement. The new-run prior predictive is a
deliberate recursion variant and should be checked against its own stated generative
convention, not assumed identical to a textbook implementation.

Phenotype compares two fixed trailing windows rather than the inferred run's
segments, without chunk restriction, coverage requirement or uncertainty. Absolute
mean and SD differences are heuristically compared; few samples default to a mean
shift. Pruning uses a value cutoff, so tied hypotheses can exceed \code{max_hyp};
the advertised hard computational cap is not strict under ties.
\repair Specify gap/new-segment semantics; verify a small enumerated Bayesian
reference for the chosen change convention. Make phenotype confidence/coverage
explicit; enforce a deterministic top-$K$ cap if that cap is required.
\checklater Quiet signal after a long gap, known mean/variance injections, tiny
reference sequences and tied hypothesis weights.

\finding{C26}{S3}{R, with source-visible approximations}{STS fidelity is not established by fitting alone}
\loc{wearable/cpd/_sts.py}{99--139 and 155--189}
STS fits a Gaussian structural likelihood, then fits t tails to standardized total
forecast errors and applies them only to simulated measurement shocks. Those errors
also contain latent-state uncertainty, so their tail weight does not uniquely
identify measurement noise. Finite standardized errors are not explicitly masked
by observed grid cells before tail fitting; the installed statsmodels treatment of
missing-cell residuals must be checked before claiming contamination. No such
library behavior was executed/assumed proven here.

Grid collisions keep the last measurement but set exogenous state indicators by
any colliding on-state row, potentially pairing an awake measurement with a sleep
indicator. Fit output counts observations before deduplication. Simulation keeps
timestamps/states but repeats draws anchored at the fitted end state on shifted
historical masks; model state/seasonal phase alignment and block independence need
library-level verification. Clipping, stationary AR(2) and independent blocks cannot
be presumed to reproduce bursts, device pile-ups or slow real variation. These are
model adequacy questions even after the coding defects are repaired.
\repair Jointly align grid values/states with one collision rule; mask tail fitting
by observed support; record fit quality and validate the full synthetic distribution,
not only fitted coefficients. Verify vendor simulation shock/anchor semantics before
claiming exact reproducibility or phase fidelity.
\checklater Colliding samples with conflicting states, long gaps, known seasonal
phase, variance/tail checks, burst autocorrelation and fresh versus repeated blocks.

\finding{C27}{S3}{D/V}{Input and setting validation does not protect numerical invariants}
\loc{wearable/steps/detect.py}{64--90 and 443--475}
\loc{wearable/cpd/drift.py}{55--59}
Settings names are checked, not numerical domains. Empty/duplicate method lists,
nonpositive/nonfinite budgets, invalid dwell times, zero stride, zero permutations,
nonpositive hazard scale, invalid direction or impossible coverage can crash or
silently misbehave. A high requested analytic drift rate may put the root outside
the fixed bracket; \code{cusum_k=0} divides by zero. Non-HRV channels without a CPD
settings block inherit HRV cadence. Existing chunk columns with some NaNs are trusted;
timestamps, strict ordering after duplicates, infinity, and array lengths are not
validated at the public detector interface. A detector lacking both calibration
hooks has no threshold entry when the pipeline invokes it.
\repair Validate domains, unique methods, channel support, timestamps, finite
raw/transform alignment and detector calibration capabilities before creating an
artifact. Reject unsupported cases with actionable errors, or implement declared
fallbacks. Do not infer that ingestibility means CPD settings are appropriate.
\checklater A compact invalid-input matrix covering each listed invariant, including
drift-only, an empty ensemble and a newly discovered detector with no score hook.

\finding{C28}{S3}{D/V}{One calibration failure can abort the entire cohort}
\loc{wearable/steps/detect.py}{398--408 and 476--495}
\code{_patient_job} does not catch per-patient failures; materializing the process
map propagates the exception and discards the usable returned results for writing.
Only the all-empty-output case marks the artifact failed. Other exceptions can leave
its manifest marked running. Global warning suppression hides numerical diagnostics.
Conversely an unarmable patient with at least two readings returns a blank output
and counts as a processed patient, without an explicit unarmable status.
\repair Preserve patient-level success/failure/unarmable reasons and a declared
cohort completeness policy; mark the artifact failed on any unhandled exception.
Scope warnings narrowly and retain diagnostics. Expose worker count rather than
always creating almost one worker per CPU for expensive STS tasks.
\checklater One patient fails while others succeed; verify manifest status, failure
table, retained outputs, and unchanged modelling population/missingness policy.

\finding{C29}{S3}{D/V}{Calibration reports can compare different computations}
\loc{scripts/cpd_calibration_compare.py}{35--47 and 55--87}
The ``latest'' selection filters ensemble, channel suffix, budget and calibration,
but not \code{on}, smoother, input artifact identity, code version or cohort/support.
It can compare STS residual detection with a newer raw shuffle run. The reader
hardcodes CSV despite configurable artifact format. Exposure is computed from
original armed rows and nominal cadence, whereas thresholds used binned support
(C20). Labels such as ``time flagged'' are row fractions, not guaranteed elapsed-time
fractions. Reported differences cannot all be attributed to calibration choice.
\repair Require a matched comparison specification and use the artifact table
reader, keyed input identity and stored exposure definition. Print mismatches rather
than selecting an arbitrary newest run.
\checklater Multiple raw/residual runs at the same budget, different smoothers and
CSV/parquet artifacts; only truly matched pairs may be compared.

\finding{C30}{S2}{V}{The present tests leave the central failure modes uncovered}
\loc{wearable/selftest.py}{543--578, 580--600 and 627--753}
The main pipeline probe uses one synthetic patient per channel and one prefix cut.
Its generated dataframe omits ingest filler rows, so C02 is not exercised. It does
not create a first-sleep-in-the-future transition. The HR check excludes the open
bin at its cut. Cuts well after a calibration date cannot reveal all future-in-day
calibration paths; alarm flags can also remain identical despite changed thresholds.
The rate bound permits 2 alarms/day against a 1/3 budget, not target agreement.
Checking a threshold's event count on the same null used to choose it is not an
independent calibration assessment. The STS coefficient/mask fixture is useful but
does not establish real-data synthetic fidelity. Existing full-path probes primarily
compare dataframe columns, not a complete as-of fit/calibration dependency ledger.
\repair Add the acceptance fixtures in this report and independent null/injection
validation; test every boundary class rather than one fixed cut. Include output
availability, fitted parameters, threshold histories, feature columns, keys and
fallback status. Preserve patient-level statistical validation separately.
\checklater Every D finding has a regression fixture; every R finding has an
explicit validation/policy decision. Run these later on an appropriate environment;
none were executed during this audit.

\finding{C31}{S4}{D/V}{Diagnostic provenance is insufficient for precise reconstruction}
\loc{wearable/steps/detect.py}{335--361 and 498--517}
The table omits most detector score streams, per-reading thresholds, per-detector
availability and the fitted hourly/state/OU parameters. Calibration rows retain
thresholds and STS parameters but not every fitted whitening state. An infinite
threshold can silently disable an uncalibratable stream without a reason code.
The row phenotype ``burn-in'' also labels settling, and drift's separate minimum-night eligibility is not represented
by the shared armed flag. Drift-only runs produce no calibration report because
\code{_thresholds} returns before adding analytic-threshold rows. The incomplete
dependency manifest (C04) makes later reconstruction still harder.
\repair Store compact versioned fit snapshots, validity/reason codes and threshold
activation intervals, with optional detailed scores. Always record analytic
thresholds and the exposure to which each rate refers. Avoid patient data in public
documentation; diagnostic artifacts remain sensitive.
\checklater A known event can be traced from raw support through fit version,
threshold, score, eligibility, publication and daily feature without rerunning an
unrecorded fitting decision.

\finding{C33}{S3}{D/V}{State inputs use an all-or-nothing availability rule}
\loc{wearable/steps/detect.py}{259--267 and 298--302}
If either sleep or steps file is absent, \code{_flags} returns None and discards the
other source too. Available sleep flags therefore do not support nightly sleep
windows when activity flags are absent; available activity is likewise ignored if
sleep is missing. No per-patient reason is emitted. An empty or poorly covered file
is a different case: state classification can become UNKNOWN, but the presence
of one sleep anywhere switches drift away from its fallback globally (C05).
\repair Load sources independently and define partially observed states explicitly.
Sleep-only information can identify sleep without asserting wake/rest/activity.
Record state coverage and fallback reasons separately from file existence.
\checklater All four sleep/steps file-presence combinations plus empty and sparse
files; usable information must be retained without inventing the absent state.

\section{Cross-cutting decisions before choosing targeted fixes or a rewrite}

\subsection{What is not itself an implementation bug}
Forward-only HMM inference with fixed parameters is valid as a specified model;
HMMs do not universally require EM training. BOCPD posterior updating is a fitting
operation even without a separate training phase. Patient-local past-only whitening
on test patients is permitted by the current pipeline and is not automatically
train/test leakage. OR fusion is a recall-oriented choice, not inherently incorrect.
STS explicitly rejects forecast surprise/tail rather than pretending to calibrate
an unexecuted forecaster. Full-day features support an end-of-day decision, not an
alert at the beginning of the same day. Each of these choices still needs accurate
claims and suitable validation.

The KCPD MMD expression uses the expected unbiased within/cross-block form; its
median kernel bandwidth is permutation-label-invariant within a fixed block.
However, exchangeability and the claimed whiteness are assumptions, not consequences
of having called a whitening function. The OU-plus-robust-update pipeline does not
establish Gaussianity, unit variance or independent tails. The Kalman pair statistic's
normal approximation likewise needs enough appropriate innovations; a small coverage
threshold is not a proof. These concerns belong in validation, not a claim that the
visible formulas are algebraically wrong.

\subsection{Minimum repair order}
\begin{enumerate}
\item Freeze the comparison contract: raw variables, transformed support, availability
timestamps, gap/reset semantics, exposure and event definitions.
\item Repair the historical cutoff, observed-row chunking, state fallback and holder
boundaries (C01, C02, C05, C07, C11, C19); certify causality only for compatible
upstream inputs (C06).
\item Repair feature coverage and dependency keys together (C03, C04). A metadata-only
fix is insufficient. Establish a new artifact-version boundary.
\item Repair null construction, missing/degenerate-score handling and repeated-event
semantics (C08--C10, C14--C18). Recalibrate after recurrence or eligibility changes.
\item Resolve refit-state coordinates and detector memory policies (C12, C13, C32), then
complete exposure, diagnostics, validation and reporting changes (C20--C31).
\item Rebuild affected CPD artifacts and dependent modelling results on the data
machine. Compare newly defined results under identical patients and feature bundles;
do not silently mix old and new artifact generations.
\end{enumerate}

\subsection{When targeted changes are sufficient}
If the intended experiment remains patient-local whitening, calibrated detectors and
daily classifier features, most D findings are localized to a small number of shared
interfaces. Preserve detector mathematics where verified; replace history slicing,
chunking, availability mapping, event representation and dependency declarations.
Changing all five detectors independently would duplicate repairs and risk changing
the experimental question. Statistical R findings still require model validation
after the plumbing is fixed.

\subsection{When a broader rewrite is justified}
A streaming or resumable service needs explicit state and clock events that the
current whole-frame API does not expose. A coherent replacement boundary would have:
\begin{enumerate}
\item A timestamp-bounded history object containing raw observations, state flags,
chunk/coverage intervals, source identity and a known availability time.
\item A fit snapshot recording normalization, OU/tail parameters and its effective
time. Refits must causally transform or reconstruct filter state.
\item One null-generation interface returning data plus the same availability,
reset and exposure definitions used by real inference.
\item Detector update interfaces with explicit memory/reset policies, score validity
and separate state-versus-verdict outputs. Repeated verdicts carry stable event IDs.
\item A publisher that never backdates information, and a daily aggregator that
preserves raw variables independently of CPD feature coverage.
\item Complete content-addressed dependencies and diagnostic fit/calibration ledgers.
\end{enumerate}
This is an architectural option, not authorization to change code in this review.

\section{Audit coverage and remaining verification}
Read in full: all five detector modules, \code{cpd/__init__.py},
\code{cpd/_whiten.py}, \code{cpd/_sts.py}, and \code{steps/detect.py}.
Inspected at their CPD boundaries: ingestion fillers, shared chunking and settings,
registry contracts, store keying, CLI construction, experiment recipes, join feature
aggregation, calibration reporting, smoke tests and the CPD selftest fixtures.
Pipeline documentation and the CPD implementation history informed interpretation.
This is not an audit of every smoother, foundation-model library, classifier or
clinical label definition. Their relevant CPD boundary risks are identified above.

No empirical sensitivity/specificity, false-alarm rate, detection delay, computational
cost or real-data impact was measured. Static review cannot establish how often the
triggering conditions occur. External library numerical behavior and STS fidelity
remain explicitly unverified. The acceptance checks define the next investigation
without requiring raw data or GPU execution on this laptop during this task.

\end{document}
